% Claim proof file for row claim_3_35 -- "LabelRoman".
% Auto-generated stub by scaffold/solve_chapter.py at row start. A manager
% agent dedicated to writing the TeX proof fills in the proof body below.
%
% This file is a sibling of `main.tex` inside `<subsection>/tex/`. The
% `subfiles` package lets it render both standalone and as part of
% `main.tex` via `\subfile{...}`.
%
% The statement is restated above the proof so the file is self-contained
% when rendered on its own. The proof must:
%   - Stay close to the lecture notes' proof if one exists (always search
%     the LN first for a `\begin{proof}` block immediately following the
%     claim; copy / lightly edit when present).
%   - Otherwise use the LN paradigm and cite prior claims by their `ref`.
%   - Be self-contained enough that a mathematician can follow it without
%     re-reading the LN.
%   - Have no `sorry`, `omitted`, or "obvious" shortcuts.

\documentclass[main]{subfiles}

\begin{document}
% ref: claim_3_35 (proof, with statement restated for readability)
% title: LabelRoman
\def\rowref{claim\_3\_35}
\phantomsection\label{claim_3_35}

% --- statement (restated) ---------------------------------------------------
\begin{Lem}[LabelRoman]\label{lem:isep_equiv_sep}
  Let $G=(J,V,E,L)$ be a CDMG.
  Define the DMG $\bar{G} := (\bar{V}, \bar{E}, \bar{L})$ with nodes $\bar{V} := V \dcup J$, directed edges $\bar{E} := E$ and bidirected edges $\bar{L} := L \cup \{j \huh j' : j, j' \in J, j \ne j'\}$.
  Then, for $A, B, C \subseteq J \cup V$, the following are equivalent:
%  \begin{enumerate}[label=(\roman*)]
%    \item $\displaystyle A \isPerp_G B \given C,$
%    \item $\displaystyle A \sPerp_{\bar{G}} B \cup J \given C,$
%    \item $\displaystyle (A\sm J) \sPerp_{\bar{G}} B \cup (J \sm C) \given C \cup (J \sm B);$
%  \end{enumerate}
  $$A \isPerp_G B \given C \iff A \sPerp_{\bar{G}} B \cup J \given C,$$
  and likewise
  $$A \idPerp_G B \given C \iff A \dPerp_{\bar{G}} B \cup J \given C.$$
\end{Lem}

% --- proof ------------------------------------------------------------------
\begin{proof}
% TODO: write the proof body.
\end{proof}

\end{document}
